Sequential robot data are signals before they become training examples. I want to understand how alignment and preprocessing alter their temporal features, how visual observations should be modeled, and which learning methods retain useful trajectory structure. My RGB-anchored CR5 pipeline and separate action-tokenization work gave me concrete data on which to ask these questions, but not yet the signal-processing depth to answer them. NTU's MSc in Signal Processing and Machine Learning would provide that depth.

For the August 2027 intake, if the current course structure remains in place, EE6401 Advanced Digital Signal Processing would deepen how I model and transform sequential data; EE6222 Machine Vision would strengthen the visual-observation layer; and EE6406 Analytic \& Ensemble Machine Learning would broaden the methods I can compare. If I choose the coursework-only option and EE6008 Collaborative Research \& Development Project is available, its three-AU, three-to-five-person team setting could vary alignment and preprocessing choices, examine their effects on temporal features, and compare how several learning methods respond. I would contribute experience with synchronized robot data and structured action representations, then carry the resulting signal-processing perspective into learning pipelines for embodied-agent R\&D.
